\documentclass[a4paper,11pt]{article}
\usepackage[T2A]{fontenc}
\usepackage[utf8]{inputenc}
\usepackage[russian]{babel}
\usepackage{amsmath,amssymb,mathtools}
\usepackage{geometry}
\usepackage{microtype}
\usepackage{tikz}
\usetikzlibrary{arrows.meta,positioning,shapes.geometric,calc}
\usepackage{booktabs,tabularx,array}
\usepackage{enumitem}
\usepackage{fancyhdr}
\usepackage{xcolor}
\usepackage{hyperref}
\usepackage{titlesec}
\usepackage{needspace}
\geometry{a4paper,left=21mm,right=21mm,top=21mm,bottom=22mm,headheight=22pt,headsep=8mm,footskip=10mm}
\definecolor{accent}{HTML}{007B7B}
\definecolor{deepaccent}{HTML}{00565A}
\definecolor{ink}{HTML}{303437}
\definecolor{softgray}{HTML}{72797D}
\color{ink}
\hypersetup{colorlinks=true,linkcolor=deepaccent,urlcolor=deepaccent}
\fancyhf{}
\fancyhead[L]{\small Факториалы и формула Лежандра}
\fancyhead[R]{\small Ксения Клыкова}
\fancyfoot[C]{\small\thepage}
\renewcommand{\headrulewidth}{0.45pt}
\renewcommand{\headrule}{\hbox to\headwidth{\color{accent}\leaders\hrule height \headrulewidth\hfill}}
\pagestyle{fancy}
\titleformat{\section}{\normalfont\Large\color{deepaccent}}{\thesection.}{0.55em}{}
\titlespacing*{\section}{0pt}{20pt}{12pt}
\titleformat{\subsection}{\normalfont\normalsize\color{deepaccent}}{}{0pt}{}
\titlespacing*{\subsection}{0pt}{12pt}{6pt}
\setlength{\parindent}{0pt}
\setlength{\parskip}{7pt}
\setlength{\emergencystretch}{1.5em}
\setlist[itemize]{leftmargin=6mm,topsep=3pt,itemsep=3pt}
\setlist[enumerate]{leftmargin=8mm,topsep=4pt,itemsep=7pt}
\newcommand{\key}[1]{\textcolor{deepaccent}{#1}}
\newcommand{\ex}[1]{\par\medskip\textcolor{deepaccent}{\textbf{Пример.} #1}\par\smallskip}
\newcommand{\nuval}[2]{v_{#1}(#2)}
\begin{document}

% ===== Титульный лист =====
\thispagestyle{empty}
\begin{center}
\vspace*{23mm}
{\Large\color{deepaccent} МАТЕМАТИКА\par}
\vspace{27mm}
{\huge Чек-лист\par}
\vspace{15mm}
{\LARGE\color{accent} Факториалы, делимость\par}
\vspace{2mm}
{\LARGE\color{accent} и формула Лежандра\par}
\vspace{20mm}
{\large Олимпиадные методы теории чисел\par}
\vfill
{\large Ксения Клыкова\par}
\vspace{12mm}
{\large 12 ноября 2026 г.\par}
\vspace{23mm}
{\small Лёвин Артём Александрович\par}
{\small эксклюзивно для Ксении Клыковой\par}
\vspace{19mm}
\begin{tikzpicture}[x=1cm,y=1cm]
\draw[accent,line width=0.8pt] (0,0) .. controls (1.6,0.7) and (3.1,-0.7) .. (4.7,0) .. controls (6.2,0.65) and (7.8,-0.65) .. (9.4,0);
\fill[accent] (4.7,0) circle (1.2pt);
\end{tikzpicture}
\end{center}
\clearpage

\section{Факториал: определение и сокращение}
Для целого $n\ge 1$ \key{факториалом} называют произведение всех натуральных чисел от $1$ до $n$:
\[
 n! = 1\cdot2\cdot3\cdots(n-1)\cdot n.
\]
По соглашению $0!=1$. Факториал определён для целых неотрицательных $n$.

\key{Главная идея:} от большого факториала удобно отделять несколько первых множителей:
\[
\boxed{n!=n(n-1)!},\qquad
\boxed{(n+1)!=(n+1)n!}.
\]
Поэтому при $n\ge1$:
\[
 \frac{(n+1)!}{n!}=n+1, \qquad
 \frac{n!}{(n-1)!}=n.
\]

\ex{Сокращение дроби с факториалами}
\[
\frac{9!}{7!}=\frac{9\cdot8\cdot7!}{7!}=9\cdot8=72.
\]
\textbf{Полезная привычка:} раскладывайте только тот факториал, у которого больше аргумент. Не вычисляйте огромные числа без необходимости.

\section{Делимость факториала на число}
Знак $a\mid b$ означает: \textit{число $b$ делится на $a$ без остатка}. Например, $72\mid 6!$.

Любое натуральное число $A>1$ раскладывается на простые множители:
\[
 A=p_1^{\alpha_1}p_2^{\alpha_2}\cdots p_r^{\alpha_r},
\]
где $p_i$ --- разные простые числа, $\alpha_i$ --- положительные целые показатели.

\key{Критерий делимости:} число $A$ делит $n!$, если в разложении $n!$ каждого простого множителя $p_i$ хватает хотя бы $\alpha_i$ раз.

\ex{Найдите наименьшее $n$, для которого $72\mid n!$}
\[
72=2^3\cdot3^2,\qquad
5!=2^3\cdot3\cdot5.
\]
В $5!$ имеется только одна тройка, поэтому $72\nmid5!$.
\[
6!=2^4\cdot3^2\cdot5\quad\Longrightarrow\quad72\mid6!.
\]
Так как все предыдущие факториалы делят $5!$, \key{наименьшее $n=6$}.

\ex{Быстрый поиск по самому крупному простому множителю}
Для $35=5\cdot7$ в $n!$ необходим множитель $7$. При $n<7$ его нет, а $7!$ содержит и $5$, и $7$. Значит, \key{$n=7$}. Аналогично, $1170=2\cdot3^2\cdot5\cdot13$ требует наличия множителя $13$, поэтому минимальное значение --- \key{$n=13$}.

\clearpage
\section{Показатель простого числа и формула Лежандра}
Обозначение $\nuval{p}{N}$ означает \key{показатель простого числа $p$} в разложении натурального числа $N$.

Например, $6!=720=2^4\cdot3^2\cdot5$, поэтому
\[
\nuval{2}{6!}=4,\qquad \nuval{3}{6!}=2,\qquad \nuval{5}{6!}=1.
\]

\key{Формула Лежандра} позволяет найти показатель простого $p$ в $n!$ без разложения всего факториала:
\[
\boxed{\displaystyle
\nuval{p}{n!}=\left\lfloor\frac np\right\rfloor+
\left\lfloor\frac n{p^2}\right\rfloor+
\left\lfloor\frac n{p^3}\right\rfloor+\cdots}
\]
Здесь $\lfloor x\rfloor$ --- округление $x$ \textit{вниз} до целого числа. Например, $\lfloor2{,}5\rfloor=2$.

\textbf{Почему формула работает?} Первый член считает числа от $1$ до $n$, кратные $p$. Второй добавляет ещё по одному множителю $p$ за каждое число, кратное $p^2$; третий --- за кратные $p^3$, и так далее.

\key{Остановка вычислений:} слагаемое становится нулём, как только $p^j>n$.

\ex{Степень двойки в $6!$}
\[
\nuval{2}{6!}=\left\lfloor\frac62\right\rfloor+
\left\lfloor\frac64\right\rfloor=3+1=\boxed{4}.
\]
Слагаемое $\lfloor 6/8\rfloor=0$ уже не требуется.

\ex{Степени двойки и тройки в $10!$}
\begin{align*}
\nuval{2}{10!}&=\left\lfloor\frac{10}{2}\right\rfloor+
\left\lfloor\frac{10}{4}\right\rfloor+
\left\lfloor\frac{10}{8}\right\rfloor=5+2+1=\boxed8,\\
\nuval{3}{10!}&=\left\lfloor\frac{10}{3}\right\rfloor+
\left\lfloor\frac{10}{9}\right\rfloor=3+1=\boxed4.
\end{align*}

\textbf{Важно:} формулу Лежандра применяют к \textit{простому} числу $p$. Для составного делителя понадобится следующий шаг.

\section{Максимальная степень числа, делящая факториал}
Пусть $A=p_1^{\alpha_1}\cdots p_r^{\alpha_r}$, где $A>1$. Ищем наибольшее целое $k\ge0$, при котором $A^k\mid n!$.

Условие делимости для каждого простого множителя:
\[
 k\alpha_i\le\nuval{p_i}{n!}.
\]
Следовательно,
\[
\boxed{\displaystyle k_{\max}=\min_{1\le i\le r}
\left\lfloor\frac{\nuval{p_i}{n!}}{\alpha_i}\right\rfloor.}
\]
\textbf{Смысл минимума:} первым закончится тот простой множитель, которого хватает на меньшее число полных наборов, необходимых для $A$.

\clearpage
\section{Как применять формулу к составным числам}
\ex{Наибольшее $k$, для которого $12^k\mid 6!$}
Число $12=2^2\cdot3$, а $6!=2^4\cdot3^2\cdot5$. Поэтому
\[
 k_{\max}=\min\left\{\left\lfloor\frac42\right\rfloor,
\left\lfloor\frac21\right\rfloor\right\}=\boxed2.
\]
Четыре двойки позволяют собрать две пары двоек; двух троек также хватает ровно на два набора.

\ex{Наибольшее $k$, для которого $18^k\mid10!$}
\[
18=2\cdot3^2,\qquad\nuval{2}{10!}=8,\qquad\nuval{3}{10!}=4.
\]
\[
 k_{\max}=\min\left\{\left\lfloor\frac81\right\rfloor,
\left\lfloor\frac42\right\rfloor\right\}=\boxed2.
\]
Ограничение создают тройки: для одной копии $18$ нужны сразу две тройки.

\ex{Олимпиадный пример: наибольшее $k$ при $12^k\mid30!$}
\begin{align*}
\nuval{2}{30!}&=\left\lfloor\frac{30}{2}\right\rfloor+
\left\lfloor\frac{30}{4}\right\rfloor+
\left\lfloor\frac{30}{8}\right\rfloor+
\left\lfloor\frac{30}{16}\right\rfloor
=15+7+3+1=26,\\
\nuval{3}{30!}&=\left\lfloor\frac{30}{3}\right\rfloor+
\left\lfloor\frac{30}{9}\right\rfloor+
\left\lfloor\frac{30}{27}\right\rfloor
=10+3+1=14.
\end{align*}
\[
 k_{\max}=\min\left\{\left\lfloor\frac{26}{2}\right\rfloor,
\left\lfloor\frac{14}{1}\right\rfloor\right\}=\boxed{13}.
\]
\textbf{Проверка максимальности:} для $k=14$ понадобилось бы $28$ двоек, но в $30!$ их только $26$.

\section{Типичные ошибки}
\begin{itemize}
\item \textbf{Путают две степени.} Величина $\nuval p{n!}$ берётся из факториала; показатель $\alpha_i$ --- из разложения делителя $A$.
\item \textbf{Забывают разделить на $\alpha_i$.} Например, при делимости на $12^k$ две двойки расходуются на каждую степень $12$.
\item \textbf{Округляют после сложения дробей.} В формуле Лежандра каждую дробь округляют вниз \textit{отдельно}.
\item \textbf{Требуют равенства показателей в самом факториале.} Лишние двойки, тройки и другие простые множители допустимы; важно, чтобы хватило всех требуемых.
\item \textbf{Берут максимум вместо минимума.} Количество полных наборов ограничивает самый дефицитный простой множитель.
\end{itemize}

% ===== Самостоятельная страница схемы =====
\clearpage
\section*{Алгоритм: максимальная степень $A$ в $n!$}
\vspace{6mm}
\begin{center}
\tikzset{
  act/.style={draw=accent,rounded corners=2.5pt,minimum width=90mm,text width=82mm,align=center,minimum height=11mm,inner sep=3.5mm,font=\small},
  decision/.style={draw=accent,diamond,aspect=2.3,text width=42mm,align=center,inner sep=1.5mm,font=\small},
  term/.style={draw=deepaccent,ellipse,minimum width=57mm,align=center,minimum height=9mm,font=\small},
  arr/.style={-{Latex[length=2.4mm]},line width=0.8pt,draw=deepaccent}
}
\begin{tikzpicture}[node distance=10mm and 14mm]
\node[term] (start) {Начало: $A>1$, $n\ge0$};
\node[act,below=of start] (factor) {Разложите $A=\prod p_i^{\alpha_i}$ на простые множители};
\node[act,below=of factor] (select) {Выберите ещё не обработанный простой множитель $p_i$};
\node[act,below=of select] (leg) {По формуле Лежандра найдите $\nuval{p_i}{n!}$};
\node[act,below=of leg] (bound) {Вычислите $b_i=\left\lfloor\nuval{p_i}{n!}/\alpha_i\right\rfloor$};
\node[decision,below=13mm of bound] (more) {Остались другие простые множители?};
\node[act,below=16mm of more] (minimum) {Выберите $k_{\max}=\min\{b_1,\ldots,b_r\}$};
\node[term,below=of minimum] (finish) {Ответ: $k_{\max}$};
\draw[arr] (start)--(factor);
\draw[arr] (factor)--(select);
\draw[arr] (select)--(leg);
\draw[arr] (leg)--(bound);
\draw[arr] (bound)--(more);
\draw[arr] (more)--node[right,font=\small,fill=white,inner sep=1.5pt]{нет}(minimum);
\draw[arr] (minimum)--(finish);
\draw[arr] (more.west)--++(-27mm,0) |- node[pos=.27,left,font=\small,fill=white,inner sep=1pt]{да} (select.west);
\end{tikzpicture}
\end{center}

% ===== Тренировка =====
\clearpage
\section*{Тренировочный блок}
\textit{Путь А. Десять задач с постепенным повышением сложности.}\par
Во всех заданиях $n$ --- целое неотрицательное число. Если требуется найти показатель делимости, указывайте наибольшее целое $k\ge0$.

\subsection{Базовый уровень}
\begin{enumerate}[label=\textcolor{deepaccent}{\arabic*.},series=tasks]
\item Вычислите $\displaystyle\frac{8!}{6!}$.
\item Упростите $\displaystyle\frac{(n+2)!}{n!}$.
\item Найдите наименьшее $n$, для которого $35\mid n!$.
\end{enumerate}
\subsection{Средний уровень}
\begin{enumerate}[label=\textcolor{deepaccent}{\arabic*.},resume=tasks]
\item Найдите наименьшее $n$, для которого $72\mid n!$.
\item По формуле Лежандра вычислите $\nuval{2}{12!}$.
\item По формуле Лежандра вычислите $\nuval{3}{25!}$.
\item Найдите наименьшее $n$, для которого $504\mid n!$.
\end{enumerate}
\subsection{Повышенный уровень}
\begin{enumerate}[label=\textcolor{deepaccent}{\arabic*.},resume=tasks]
\item Найдите наибольшее $k$, при котором $12^k\mid15!$.
\item Найдите наибольшее $k$, при котором $18^k\mid20!$.
\item Найдите наибольшее $k$, при котором $72^k\mid30!$.
\end{enumerate}

% ===== Ответы =====
\clearpage
\section*{Ответы к тренировочному блоку}
\vspace{6mm}
\renewcommand{\arraystretch}{1.65}
\begin{center}
\begin{tabularx}{\linewidth}{@{}>{\centering\arraybackslash}p{16mm} X@{}}
\toprule
\textbf{№} & \textbf{Ответ}\\
\midrule
1 & $56$.\\
2 & $(n+2)(n+1)$.\\
3 & $n=7$.\\
4 & $n=6$.\\
5 & $10$.\\
6 & $10$.\\
7 & $n=7$.\\
8 & $k=5$.\\
9 & $k=4$.\\
10 & $k=7$.\\
\bottomrule
\end{tabularx}
\end{center}
\vspace{8mm}
\subsection*{Числовые ориентиры для самопроверки}
В задачах 8--10 полезно сверить показатели простых множителей, прежде чем брать минимум:
\begin{align*}
\text{№ 8:}\quad&\nuval{2}{15!}=11,\quad \nuval{3}{15!}=6;\\
\text{№ 9:}\quad&\nuval{2}{20!}=18,\quad \nuval{3}{20!}=8;\\
\text{№ 10:}\quad&\nuval{2}{30!}=26,\quad \nuval{3}{30!}=14.
\end{align*}

\vfill
\begin{center}\small\color{softgray}
Лёвин Артём Александрович\par
эксклюзивно для Ксении Клыковой
\end{center}
\end{document}
